\documentclass[../Main.tex]{subfiles}


\begin{document}

\section{Reflexion}

\subsection{Reflexion des Arbeitsprozesses und der Gruppenarbeit}

Unser Projekt sowie unsere Zusammenarbeit war von deutlichen Fortschritten, aber auch von realen Schwierigkeiten geprägt.

In der Zusammenarbeit haben sich unsere unterschiedlichen Fähigkeiten als grosser Vorteil erwiesen. Wenn einer von uns irgendwo nicht weiterkam, konnte der andere oft aushelfen, und die Aufteilung der Aufgaben funktionierte meist reibungslos. Unsere Kommunikation, die bereits auf einem guten Niveau war, hat sich im Verlauf des Projekts weiter verbessert, und wir haben uns gegenseitig angespornt, was das Projekt massgeblich vorangetrieben hat. Die grösste Schwierigkeit war hingegen die zeitliche Koordination. Da wir zu unterschiedlichen Zeiten unter der Woche arbeiten, entstand bei gegenseitigen Abhängigkeiten oft ein Bottle-Neck. Hinzu kommt, dass wir nicht immer dieselbe Vision für das Spiel hatten und erst dann weiterarbeiteten, wenn wir einen Kompromiss gefunden hatten. Das hat zwar Zeit gekostet, war aber wichtig, da die Arbeit nur dann bestmöglich bewältigt werden kann, wenn beide mit voller Überzeugung hinter dem Projekt stehen.

Auf technischer Ebene übertrifft die Qualität der bisher umgesetzten Systeme das, was wir ursprünglich erwartet hatten. Besonders hervorzuheben sind die Flexibilität der Frame Data, der FloatingBody-Algorithmus, der eine geschmeidige Bewegung auf unebenen Oberflächen ermöglicht, sowie die State Machine für den Spieler-Controller. Bewusst haben wir dabei zunächst eine solide Basis aufgebaut, bevor wir die darauf aufbauenden Systeme implementieren. Das hatte jedoch zur Folge, dass uns noch einige Features fehlen. Bei der Assetproduktion lief es weniger gut, da die Lernkurve steil ist und wir noch nicht genug Zeit hatten, uns vollständig einzuarbeiten. Das Animationsproblem haben wir teilweise gelöst, indem wir ein Assetpack erworben und die enthaltenen Animationen anschliessend angepasst haben.

\subsection{Plan für die Weiterentwicklung}

Foobar

\end{document}
